% ===========================================================================
% Poste 11 — Montant pour le soutien des aînés
% ===========================================================================
\poste{11}{Montant pour le soutien des aînés}{QC\_aines}

\pdesc Crédit d'impôt remboursable québécois pour les personnes de 70~ans et
plus à revenu faible ou moyen, réclamé à la ligne 463 de la déclaration.

\pobj Soutenir le revenu des aînés les plus âgés et les plus vulnérables,
au-delà des pensions fédérales.

\pregle Chaque adulte de 70~ans et plus ouvre droit à un montant maximal $M$
--- un couple dont les deux conjoints sont admissibles part donc de $2M$. Le
total est réduit selon le revenu familial net au-delà du seuil $s$ (qui
dépend du nombre d'adultes, pas du nombre d'aînés) :
\[
\text{crédit} = \pos{\,k \times M - \tau \times \pos{\RFN - s}\,}
\qquad k \in \{1, 2\}
\]
où $k$ est le nombre d'adultes de 70~ans et plus. Si aucun adulte n'a
70~ans, le crédit est nul.

\pparams

\begin{center}
\small
\begin{tabular}{l r r}
\toprule
Paramètre & 2025 & 2026 \\
\midrule
Montant par aîné admissible ($M$)   & \D{2000}  & \D{2000} \\
Âge d'admissibilité                 & 70 ans    & 70 ans \\
Seuil de réduction --- 1 adulte ($s$)  & \D{27835} & \D{28405} \\
Seuil de réduction --- couple ($s$)    & \D{45270} & \D{46200} \\
Taux de réduction ($\tau$)          & \pc{5.4}  & \pc{5.47} \\
\bottomrule
\end{tabular}
\end{center}

% ============================================================================
% GÉNÉRÉ par scripts/exemples-guide.ts — NE PAS ÉDITER À LA MAIN.
% Régénérer : npm run exemples (chaque chiffre est vérifié contre le moteur).
% ============================================================================
\pexemples \D{2000} par adulte de 70~ans et plus,
réduits de \pc{5.4} du $\RFN$ excédant le seuil (\D{27835}
pour un adulte seul ; \D{45270} pour un couple).

\medskip\noindent\textbf{M10 --- retraité vivant seul, 70~ans, revenu de pension privée de \D{20000}.}
\begin{align*}
\text{réduction} &= \pos{\num{29891.99} - \num{27835}} \times \pc{5.4} = \num{111.08} \\
\text{crédit} &= \num{2000} - \num{111.08} = \num{1888.92}
\end{align*}
Le $\RFN$ (\num{29891.99}) inclut la Sécurité de la vieillesse, pas seulement
la pension privée. Crédit : \D{1888.92}.

\medskip\noindent\textbf{M11 --- couple de retraités, 72 et 70~ans, pensions privées de \D{30000} et \D{10000}.}
Les deux conjoints ont 70~ans et plus : montant maximal de
$2 \times \num{2000} = \num{4000}$.
\begin{align*}
\text{réduction} &= \pos{\num{57582.28} - \num{45270}} \times \pc{5.4} = \num{664.86} \\
\text{crédit} &= \num{4000} - \num{664.86} = \num{3335.14}
\end{align*}
Crédit : \D{3335.14}.

\medskip\noindent\textbf{M12 --- retraité vivant seul, 76~ans, revenu de pension privée de \D{110000}.}
Le $\RFN$ (\num{115737.85}) est si élevé que la réduction
(\num{4746.75}) dépasse le montant maximal : crédit \emph{éteint}
(\D{0}). À noter : M13 (68 et 66~ans) n'y a pas droit non plus —
l'âge d'admissibilité est 70~ans, pas 65.


\prefs Loi sur les impôts (RLRQ, c.~I-3) ; déclaration TP-1, ligne 463 ;
Revenu Québec ; ministère des Finances du Québec, \emph{Paramètres du
régime d'imposition}.
